% 对应题号 M2-104，本题条件、变量及几何关系逐项核对。
% 修订：启用中文断行；核对题设及图文对应。
\documentclass[tikz,border=14pt]{standalone}
\usepackage{fontspec,xeCJK}
\usepackage{amsmath,amssymb}
\usepackage{xcolor}
\usepackage{tikz}
\usetikzlibrary{arrows.meta,calc,patterns,positioning,bending}
\IfFontExistsTF{Heiti SC}{\IfFontExistsTF{Heiti SC}{\setmainfont{Heiti SC}}{\setmainfont{FandolHei-Regular}}\IfFontExistsTF{Heiti SC}{\setsansfont{Heiti SC}}{\setsansfont{TeX Gyre Heros}}}{\IfFontExistsTF{Songti SC}{\setmainfont{Songti SC}}{\setmainfont{FandolSong-Regular}}\IfFontExistsTF{Songti SC}{\setsansfont{Songti SC}}{\setsansfont{FandolSong-Regular}}}
\definecolor{bluefield}{RGB}{24,76,255}
\definecolor{redfield}{RGB}{232,49,35}
\definecolor{orangefield}{RGB}{236,132,18}
\definecolor{greenfield}{RGB}{18,132,86}
\definecolor{purplefield}{RGB}{112,72,190}
\definecolor{linegray}{RGB}{60,68,78}
\definecolor{softblue}{RGB}{238,243,255}
\definecolor{softgreen}{RGB}{238,249,244}
\definecolor{softorange}{RGB}{255,247,235}
\tikzset{
  arrow/.style={-{Latex[length=3.1mm,width=2.1mm]},line width=1.1pt},
  thinarr/.style={-{Latex[length=2.5mm,width=1.7mm]},line width=.85pt},
  axis/.style={-{Latex[length=2.5mm,width=1.7mm]},line width=.8pt,draw=linegray},
  guide/.style={draw=linegray!55,dashed,line width=.65pt},
  box/.style={draw=linegray,rounded corners=3pt,fill=white,align=center,inner sep=6pt,line width=.75pt},
  bluebox/.style={box,draw=bluefield,fill=softblue},
  greenbox/.style={box,draw=greenfield,fill=softgreen},
  orangebox/.style={box,draw=orangefield,fill=softorange}
}
\IfFontExistsTF{Heiti SC}{\setCJKmainfont[AutoFakeBold=1.5]{Heiti SC}}{\setCJKmainfont[AutoFakeBold=1.5]{FandolHei-Regular}}
\IfFontExistsTF{Heiti SC}{\setCJKsansfont{Heiti SC}}{\setCJKsansfont{FandolHei-Regular}}
\begin{document}

\begin{tikzpicture}[x=1cm,y=1cm]
\path[use as bounding box] (-10,-4.6) rectangle (10,4.6);
\fill[white] (-10,-4.6) rectangle (10,4.6);
\node[font=\Large\bfseries,text width=18.6cm,align=center] at (0,4.05) {半无限条带的单模边界：横向正弦模态沿纵向指数衰减};
\begin{scope}[shift={(-4.8,-.1)}]
\fill[softblue] (0,0) rectangle (4.2,3); \draw[linegray,line width=1pt] (0,0)--(4.2,0)--(4.2,3); \draw[linegray,line width=1pt] (0,0)--(0,3);
\draw[redfield,line width=1.3pt] plot[domain=0:4.2,samples=70] (\x,{.65*sin(180*\x/4.2)});
\foreach \y in {.7,1.4,2.1,2.8}{\draw[bluefield!70,line width=.8pt] plot[domain=0:4.2,samples=60] (\x,{\y+.45*exp(-\y)*sin(180*\x/4.2)});}
\node[below] at (2.1,0) {$0<x<a$}; \node[left] at (0,1.5) {$y>0$};
\node[font=\small,align=center,text width=5.5cm]at(2.15,-2){曲线为不同高度的电势剖面示意};
\end{scope}
\begin{scope}[shift={(4.65,-.1)}]
\node[bluebox,text width=6.9cm] at (0,1.65) {$X''+k^2X=0,\quad X(0)=X(a)=0\Rightarrow k_n=n\pi/a$};
\node[greenbox,text width=6.9cm] at (0,0) {边界只有 $n=1$ 模态，因此无需叠加其余本征函数。};
\node[orangebox,text width=6.9cm] at (0,-1.7) {$\displaystyle \varphi(x,y)=V_0\sin\frac{\pi x}{a}\,e^{-\pi y/a}$};
\end{scope}
\end{tikzpicture}
\end{document}
